%% kapa: exam=2 round=35 year=2024 period=1 subject=appraisal_practice subject_ko=감정평가실무 seq=1 number=1 points=40
%% source: Q-Net 감정평가사 기출문제 (cst00309) artlSeq=5213532 file=raw/1008/5213532/2024년 제35회 감정평가사 2차시험 문제지.pdf
%% conversion: kordoc 4.24.51 -> tools/gen2.py (auto); review=proofread-sonnet; flags=none
\begin{longquestion}{【문제 1】}
감정평가사 甲은 중앙토지수용위원회로부터 수용재결평가를 의뢰받았다. 관련 법규 및 이론을 참고하고, 제시된 자료를 활용하여 다음의 물음에 답하시오. (40점) \par
\medskip\noindent 물음 1) 대상토지의 감정평가에서 < 자료 4 >의 연도별 표준지공시지가 중 적정한 비교표준지의 공시기준일을 정하고 그 근거를 제시하시오. (10점) \par
\medskip\noindent 물음 2) 대상토지의 감정평가에서 적용할 지가변동률은 첫째, 비교표준지가 소재하는 해당 시·군·구의 용도지역별 지가변동률, 둘째, 해당 공익사업과 관계없는 인근 시·군·구의 용도지역별 지가변동률 중 어느 것인지를 결정하고 그 근거를 제시하시오. (10점) \par
\medskip\noindent 물음 3) 대상토지에 대하여 시산가액에 대한 합리성 검토는 생략하고 보상평가액을 산정하시오. (10점) \par
\medskip\noindent 물음 4) 대상지장물인 건물에 대한 보상평가액을 < 자료 10 >과 < 자료 11 >을 참고하여 산정하시오. (10점) \par
\medskip\noindent < 자료 1 > 평가개요 \par
1. 의뢰인: 중앙토지수용위원회위원장 \par
2. 사업시행자: ○○공사 \par
3. 평가목적: 수용재결 \par
4. 평가의뢰일자: 2024. 6. 1. \par
\hspace*{1.2em}제시된 가격시점(재결일): 2024. 7. 1. \par
\hspace*{1.2em}사업인정고시일: 2015. 12. 30. \par
\medskip\noindent < 자료 2 > 사업개요 \par
1. 사업명: ○○산업단지 재생사업지구 기반시설공사 \par
2. 사업시행지의 위치: A광역시 B구 C동 10번지 일원 \par
3. 사업의 규모: 500,000m\textsuperscript{2} \par
4. 사업추진 경위 \par
\hspace*{1em}- 2009. 9. 30.: 재생사업 우선사업지구 선정(국토교통부) \par
\hspace*{1em}- 2010. 2. 25.: ○○산업단지 재생계획 수립용역 착수 \par
\hspace*{1em}- 2013. 12. 30.: 재생사업지구 지정(재생계획) 및 지형도면 고시 \par
\hspace*{1em}- 2015. 12. 30.: 재생사업지구 지정(재생계획) 변경, 재생시행계획 승인 고시 및 지형도면 고시 \par
\medskip\noindent < 자료 3 > 의뢰목록 \par
1. 대상토지 목록 \par
{\small\centering
\resizebox{\linewidth}{!}{%
\begin{tabular}{|c|c|c|c|c|c|c|c|c|c|c|}
\hline
\makecell{일련\\번호} & 소재지 & 지목 & \makecell{이용\\상황} & \makecell{전체\\면적\\(m\textsuperscript{2})} & \makecell{편입\\면적\\(m\textsuperscript{2})} & \makecell{용도\\지역} & \makecell{도로\\교통} & \makecell{형상\\지세} & \makecell{주위\\환경} & \makecell{피수\\용자} \\
\hline
1 & \makecell{B구\\C동 10} & \makecell{공장\\용지} & 공업용 & 990 & 990 & \makecell{일반\\공업} & \makecell{소로\\한면} & \makecell{사다리형\\평지} & \makecell{기존\\공장\\지대} & 乙 \\
\hline
\end{tabular}}\par}

2. 지장물 목록 \par
{\small\centering
\resizebox{\linewidth}{!}{%
\begin{tabular}{|c|c|c|c|c|c|c|c|c|}
\hline
\makecell{일련\\번호} & 소재지 & \makecell{물건의\\종류} & 구조 및 규격 & \makecell{수량\\(m\textsuperscript{2})} & 사용승인일 & 등급 & \makecell{내용\\년수\\(년)} & \makecell{피수\\용자} \\
\hline
가 & \makecell{B구\\C동 10} & 공장 & \makecell{철골조\\철골지붕틀\\샌드위치판넬\\잇기} & 660 & 2012. 3. 2. & 중급 & 40 & 乙 \\
\hline
\end{tabular}}\par}

\medskip\noindent < 자료 4 > 사업지구내 표준지공시지가 현황 \par
{\small\centering
\resizebox{\linewidth}{!}{%
\begin{tabular}{|c|c|c|c|c|c|c|}
\hline
\multicolumn{2}{|c|}{기호} & 1 & 2 & 3 & 4 & 5 \\
\hline
\multicolumn{2}{|c|}{소재지} & B구 C동 70 & B구 C동 100 & B구 C동 200 & B구 C동 300 & B구 C동 400 \\
\hline
\multicolumn{2}{|c|}{면적(m\textsuperscript{2})} & 850 & 450 & 600 & 290 & 1,800 \\
\hline
\multicolumn{2}{|c|}{지목} & 공장용지 & 대 & 공장용지 & 공장용지 & 대 \\
\hline
\multicolumn{2}{|c|}{이용상황} & 공업용 & 주상용 & 공업용 & 공업용 & 공업용 \\
\hline
\multicolumn{2}{|c|}{용도지역} & 일반공업 & 일반공업 & 일반공업 & 일반공업 & 일반공업 \\
\hline
\multicolumn{2}{|c|}{도로교통} & 소로각지 & 중로한면 & 소로한면 & 세로(가) & 중로한면 \\
\hline
\multicolumn{2}{|c|}{형상지세} & 세장형 평지 & 세장형 평지 & 부정형 평지 & 세장형 평지 & 가장형 평지 \\
\hline
\multirow{4}{*}{\makecell{연도별\\표준지\\공시지가\\(원/m\textsuperscript{2})}} & 2013.1.1 & 700,000 & 900,000 & 690,000 & 610,000 & 850,000 \\
\cline{2-7}
 & 2014.1.1 & 770,000 & 970,000 & 760,000 & 680,000 & 920,000 \\
\cline{2-7}
 & 2015.1.1 & 850,000 & 1,030,000 & 840,000 & 750,000 & 990,000 \\
\cline{2-7}
 & 2024.1.1 & 1,230,000 & 1,550,000 & 1,200,000 & 1,000,000 & 1,500,000 \\
\hline
\end{tabular}}\par}

주1) 주위환경은 공히 기존공장지대임 \par
주2) 대상토지와 지리적으로 인접한 정도는 기호 1, 2, 3, 4, 5번의 순서임 \par
\medskip\noindent < 자료 5 > A광역시 B구 전체 표준지공시지가의 평균변동률 \par
{\small\centering
\begin{tabular}{|c|c|c|c|c|}
\hline
\multirow{2}{*}{구분} & 2014년/2013년 & 2015년/2013년 & 2024년/2013년 & 2024년/2015년 \\
\cline{2-5}
 & 3.523\% & 7.179\% & 12.055\% & 11.412\% \\
\hline
\end{tabular}\par}

\medskip\noindent < 자료 6 > 지가변동률 현황 \par
- A광역시 \par
{\small\centering
\begin{tabular}{|c|c|c|}
\hline
구분 & 기간 & 지가변동률(\%) \\
\hline
\multirow{2}{*}{A광역시 공업지역} & 2013. 12. 30. \textasciitilde{} 2024. 7. 1. & 12.321(1.12321) \\
\cline{2-3}
 & 2015. 12. 30. \textasciitilde{} 2024. 7. 1. & 10.850(1.10850) \\
\hline
\end{tabular}\par}

- A광역시 B구 \par
{\small\centering
\begin{tabular}{|c|c|c|}
\hline
구분 & 기간 & 지가변동률(\%) \\
\hline
\multirow{5}{*}{B구 공업지역} & 2013. 1. 1. \textasciitilde{} 2013. 12. 29. & 3.795(1.03795) \\
\cline{2-3}
 & 2013. 12. 30. \textasciitilde{} 2024. 7. 1. & 36.158(1.36158) \\
\cline{2-3}
 & 2015. 1. 1. \textasciitilde{} 2015. 12. 29. & 3.662(1.03662) \\
\cline{2-3}
 & 2015. 12. 30. \textasciitilde{} 2024. 7. 1. & 19.450(1.19450) \\
\cline{2-3}
 & 2024. 1. 1. \textasciitilde{} 2024. 7. 1. & 0.225(1.00225) \\
\hline
\end{tabular}\par}

- A광역시 B구와 인접한 인근 시·군·구 \par
{\small\centering
\begin{tabular}{|c|c|c|}
\hline
구분 & 기간 & 지가변동률(\%) \\
\hline
\multirow{3}{*}{C구 공업지역} & 2013. 12. 30. \textasciitilde{} 2024. 7. 1. & 29.092(1.29092) \\
\cline{2-3}
 & 2015. 12. 30. \textasciitilde{} 2024. 7. 1. & 25.350(1.25350) \\
\cline{2-3}
 & 2024. 1. 1. \textasciitilde{} 2024. 7. 1. & 2.358(1.02358) \\
\hline
\end{tabular}\par}

{\small\centering
\begin{tabular}{|c|c|c|}
\hline
구분 & 기간 & 지가변동률(\%) \\
\hline
\multirow{3}{*}{D구 공업지역} & 2013. 12. 30. \textasciitilde{} 2024. 7. 1. & 15.355(1.15355) \\
\cline{2-3}
 & 2015. 12. 30. \textasciitilde{} 2024. 7. 1. & 13.270(1.13270) \\
\cline{2-3}
 & 2024. 1. 1. \textasciitilde{} 2024. 7. 1. & 2.032(1.02032) \\
\hline
\end{tabular}\par}

{\small\centering
\begin{tabular}{|c|c|c|}
\hline
구분 & 기간 & 지가변동률(\%) \\
\hline
\multirow{3}{*}{E구 공업지역} & 2013. 12. 30. \textasciitilde{} 2024. 7. 1. & 17.266(1.17266) \\
\cline{2-3}
 & 2015. 12. 30. \textasciitilde{} 2024. 7. 1. & 15.850(1.15850) \\
\cline{2-3}
 & 2024. 1. 1. \textasciitilde{} 2024. 7. 1. & 1.985(1.01985) \\
\hline
\end{tabular}\par}

주1) 해당월의 지가변동률이 미고시된 경우에는 직전월의 지가변동률을 연장하여 적용함 \par
주2) 생산자물가지수는 고려하지 않기로 함 \par
\medskip\noindent < 자료 7 > 거래사례 현황 \par
{\small\centering
\resizebox{\linewidth}{!}{%
\begin{tabular}{|c|c|c|c|c|c|c|c|c|c|}
\hline
기호 & 소재지 & \makecell{면적\\(m\textsuperscript{2})} & 지목 & \makecell{이용\\상황} & \makecell{용도\\지역} & \makecell{도로\\교통} & \makecell{형상\\지세} & \makecell{거래단가\\(원/m\textsuperscript{2})} & 거래시점 \\
\hline
가 & \makecell{B구\\C동 30} & 780 & \makecell{공장\\용지} & 공업용 & \makecell{일반\\공업} & \makecell{소로\\한면} & \makecell{사다리\\평지} & 1,600,000 & 2013. 2. 18. \\
\hline
나 & \makecell{B구\\D동 55} & 950 & \makecell{공장\\용지} & 공업용 & \makecell{일반\\공업} & \makecell{소로\\각지} & \makecell{세장형\\평지} & 2,400,000 & 2015. 5. 6. \\
\hline
\end{tabular}}\par}

주1) 기호 가)는 인근지역의 기존공장지대이고, 매수자가 양도소득세를 부담하는 조건으로 거래된 사례로 조사되며, 2013. 2. 18. \textasciitilde{} 2024. 7. 1.의 공업지역 지가변동률은 38.456\%임 \par
주2) 기호 나)는 인근지역의 기존공장지대이고, 특수관계인간의 거래로서 고가로 거래된 사례로 조사되며, 2015. 5. 6. \textasciitilde{} 2024. 7. 1.의 공업지역 지가변동률은 20.795\%임 \par
\medskip\noindent < 자료 8 > 평가사례 현황 \par
{\small\centering
\resizebox{\linewidth}{!}{%
\begin{tabular}{|c|c|c|c|c|c|c|c|c|c|c|}
\hline
\makecell{기\\호} & 소재지 & \makecell{면적\\(m\textsuperscript{2})} & 지목 & \makecell{이용\\상황} & \makecell{용도\\지역} & \makecell{도로\\교통} & \makecell{형상\\지세} & \makecell{보상가액\\(원/m\textsuperscript{2})} & \makecell{가격\\시점} & \makecell{평가\\목적} \\
\hline
다 & \makecell{B구\\C동 330} & 880 & \makecell{공장\\용지} & 공업용 & \makecell{일반\\공업} & \makecell{소로\\한면} & \makecell{사다리\\평지} & 2,600,000 & \makecell{2021.\\12. 25.} & \makecell{협의\\보상} \\
\hline
라 & \makecell{D구\\F동 100} & 1,100 & \makecell{공장\\용지} & 공업용 & \makecell{일반\\공업} & \makecell{소로\\한면} & \makecell{세장형\\평지} & 1,400,000 & \makecell{2022.\\5. 30.} & \makecell{협의\\보상} \\
\hline
마 & \makecell{E구\\G동 180} & 1,050 & \makecell{공장\\용지} & 공업용 & \makecell{일반\\공업} & \makecell{소로\\한면} & \makecell{사다리\\평지} & 1,800,000 & \makecell{2022.\\7. 25.} & 담보 \\
\hline
\end{tabular}}\par}

주1) 기호 다)는 인근지역의 기존공장지대로 ○○산업단지 재생사업지구 기반시설공사에 편입되어 협의보상이 완료된 사례이고, 2021. 12. 25. \textasciitilde{} 2024. 7. 1.의 공업지역 지가변동률은 3.456\%임 \par
주2) 기호 라)는 동일수급권 유사지역의 기존공장지대로 ○○관리소 건설공사에 편입되어 협의보상이 완료된 사례이고, 해당 공익사업의 시행에 따른 가격의 변동이 반영되어 있지 않다고 인정되며, 2022. 5. 30. \textasciitilde{} 2024. 7. 1.의 공업지역 지가변동률은 2.495\%임 \par
주3) 기호 마)는 동일수급권 유사지역의 기존공장지대에 위치하며, 담보목적의 평가사례로서 2022. 7. 25. \textasciitilde{} 2024. 7. 1.의 공업지역 지가변동률은 2.793\%임 \par
\medskip\noindent < 자료 9 > 가치형성요인 비교자료 \par
- 지역요인 격차율 \par
비교표준지가 있는 지역의 표준적인 획지의 최유효이용과 사례가 있는 지역의 표준적인 획지의 최유효이용을 판정하여 비교함. B구 C동(비교표준지 소재)은 D구 F동(사례 소재)과 지역요인을 비교하였을 때 가로의 계통 및 연속성 등에서 5\%정도 우세하며, 산업도로 등과의 접근의 용이성에서 20\%정도 우세한 것으로 나타남. 그 외 지역 간 지역요인은 상호 대등한 것으로 판단됨 \par
- 개별요인 격차율 \par
{\small\centering
\begin{tabular}{|c|c|c|c|c|c|}
\hline
\multirow{5}{*}{토지용도} &  & 주거용 & 상업업무 & 주상복합 & 공업용 \\
\cline{2-6}
 & 주거용 & 1.00 & 1.43 & 1.20 & 0.98 \\
\cline{2-6}
 & 상업업무 & 0.70 & 1.00 & 0.84 & 0.69 \\
\cline{2-6}
 & 주상복합 & 0.83 & 1.19 & 1.00 & 0.82 \\
\cline{2-6}
 & 공업용 & 1.02 & 1.46 & 1.22 & 1.00 \\
\hline
\end{tabular}\par}

{\small\centering
\begin{tabular}{|c|c|c|c|c|c|}
\hline
\multirow{5}{*}{형상} &  & 정방형 & \makecell{(가로·세로)\\장방형} & 사다리형 & 부정형 \\
\cline{2-6}
 & 정방형 & 1.00 & 1.00 & 0.98 & 0.95 \\
\cline{2-6}
 & \makecell{(가로·세로)\\장방형} & 1.00 & 1.00 & 0.98 & 0.95 \\
\cline{2-6}
 & 사다리형 & 1.02 & 1.02 & 1.00 & 0.97 \\
\cline{2-6}
 & 부정형 & 1.05 & 1.05 & 1.03 & 1.00 \\
\hline
\end{tabular}\par}

{\small\centering
\begin{tabular}{|c|c|c|c|c|c|c|c|}
\hline
\multirow{7}{*}{\makecell{도로\\접면}} &  & 중로한면 & 중로각지 & 소로한면 & 소로각지 & 세로(가) & 세각(가) \\
\cline{2-8}
 & 중로한면 & 1.00 & 1.06 & 0.89 & 0.94 & 0.84 & 0.88 \\
\cline{2-8}
 & 중로각지 & 0.94 & 1.00 & 0.84 & 0.88 & 0.79 & 0.83 \\
\cline{2-8}
 & 소로한면 & 1.12 & 1.19 & 1.00 & 1.05 & 0.94 & 0.99 \\
\cline{2-8}
 & 소로각지 & 1.07 & 1.14 & 0.95 & 1.00 & 0.90 & 0.94 \\
\cline{2-8}
 & 세로(가) & 1.19 & 1.27 & 1.06 & 1.12 & 1.00 & 1.05 \\
\cline{2-8}
 & 세각(가) & 1.13 & 1.21 & 1.01 & 1.06 & 0.95 & 1.00 \\
\hline
\end{tabular}\par}

주) 격차율 자료를 이용한 요인 비교치 산정은 상승식을 적용할 것 \par
\medskip\noindent < 자료 10 > 건물 재조달원가 자료 \par
- 재생사업지구 인근에 대상건물과 유사하게 신축예정인 공장건물이 있다. 아래의 “공장건물 신축개요”와 “자본환원표”를 참고하여 산정한 재조달원가를 대상건물의 재조달원가로 적용할 것 \par
- 공장건물 신축개요 \par
{\small\centering
\begin{tabular}{|c|c|c|c|}
\hline
\multicolumn{2}{|c|}{구 분} & 내 용 & 비 고 \\
\hline
\multicolumn{2}{|c|}{도급금액} & 630,000,000원 &  \\
\hline
\multicolumn{2}{|c|}{건축연면적} & 700m\textsuperscript{2} &  \\
\hline
\multicolumn{2}{|c|}{건축구조} & \makecell{철골조 철골지붕틀\\샌드위치판넬잇기} &  \\
\hline
\multicolumn{2}{|c|}{\makecell{건축공사계약일\\및 공사기간}} & 2024. 7. 1.부터 1년간 &  \\
\hline
\multirow{3}{*}{\makecell{도급\\금액\\지급\\일정}} & 계약금 & 계약시 도급금액의 10\% & \multirow{3}{*}{\makecell{소유자는 잔금을\\준공시점에 건물을 담보로\\대출 받아 지급할 예정임}} \\
\cline{2-3}
 & 중도금 & 6개월 후 도급금액의 20\% &  \\
\cline{2-3}
 & 잔금 & 준공시점에 도급금액의 70\% &  \\
\hline
\multicolumn{2}{|c|}{대출조건} & \makecell{대출이율 연 5.0\%,\\10년간 매월 원리금균등분할상환} &  \\
\hline
\multicolumn{2}{|c|}{시장조건} & 시장이자율 연 6.0\% &  \\
\hline
\multicolumn{2}{|c|}{상환계획} & \multicolumn{2}{c|}{\makecell{대출을 받고 4년 후의 시점에\\미상환잔액을 일시에 상환하려고 함}} \\
\hline
\end{tabular}\par}

- 자본환원표 \par
□ 이자율 연 5.0\% \par
{\small\centering
\resizebox{\linewidth}{!}{%
\begin{tabular}{|c|c|c|c|c|c|c|}
\hline
기간(월) & \makecell{일시불\\내가계수} & \makecell{연금\\내가계수} & \makecell{감채기금\\계수} & \makecell{일시불\\현가계수} & \makecell{연금\\현가계수} & 저당상수 \\
\hline
6개월 & 1.025262 & 6.062848 & 0.164939 & 0.975361 & 5.913463 & 0.169106 \\
\hline
12개월 & 1.051162 & 12.278855 & 0.081441 & 0.951328 & 11.681222 & 0.085607 \\
\hline
48개월 & 1.220895 & 53.014885 & 0.018863 & 0.819071 & 43.422956 & 0.023029 \\
\hline
60개월 & 1.283359 & 68.006083 & 0.014705 & 0.779205 & 52.990706 & 0.018871 \\
\hline
72개월 & 1.349018 & 83.764259 & 0.011938 & 0.741280 & 62.092777 & 0.016105 \\
\hline
120개월 & 1.647009 & 155.282279 & 0.006440 & 0.607161 & 94.281350 & 0.010607 \\
\hline
\end{tabular}}\par}

□ 이자율 연 6.0\% \par
{\small\centering
\resizebox{\linewidth}{!}{%
\begin{tabular}{|c|c|c|c|c|c|c|}
\hline
기간(월) & \makecell{일시불\\내가계수} & \makecell{연금\\내가계수} & \makecell{감채기금\\계수} & \makecell{일시불\\현가계수} & \makecell{연금\\현가계수} & 저당상수 \\
\hline
6개월 & 1.030378 & 6.075502 & 0.164595 & 0.970518 & 5.896384 & 0.169595 \\
\hline
12개월 & 1.061678 & 12.335562 & 0.081066 & 0.941905 & 11.618932 & 0.086066 \\
\hline
48개월 & 1.270489 & 54.097832 & 0.018485 & 0.787098 & 42.580318 & 0.023485 \\
\hline
60개월 & 1.348850 & 69.770031 & 0.014333 & 0.741372 & 51.725561 & 0.019333 \\
\hline
72개월 & 1.432044 & 86.408856 & 0.011573 & 0.698302 & 60.339514 & 0.016573 \\
\hline
120개월 & 1.819397 & 163.879347 & 0.006102 & 0.549633 & 90.073453 & 0.011102 \\
\hline
\end{tabular}}\par}

\medskip\noindent < 자료 11 > 건물의 감가수정 자료 \par
- 대상건물의 감가수정액은 아래의 내용을 모두 참작하여 적용할 것 \par
- 대상건물에 대하여 시간의 경과, 노후화 등에 따른 물리적 감가수정은 정액법을 적용하고, 경과연수는 연단위로 산정하며, 잔가율 등은 고려하지 아니함 \par
- 대상건물은 건축당시 층고가 낮게 설계되어 정상적인 임대료 대비 월임대료는 50원/m\textsuperscript{2} 낮은 실정이다. 이러한 문제에 관하여 조임대료승수(GRM: Gross Rent Multiplier)를 활용하여 감가수정액을 산정하며, GRM은 12를 적용하기로 함 \par
- 대상건물은 인접한 혐오시설에 기인한 악취 등으로 인하여 정상적인 임대료 대비 월임대료는 100원/m\textsuperscript{2} 낮은 실정이다. 이 문제로 인한 임대료손실액 중 건물부분이 20\%를 차지하는 것으로 판단된다. 환원율은 연 6\%를 적용하기로 함 \par
\medskip\noindent < 자료 12 > 기타 사항 \par
1. 공시지가변동률 산정시 백분율로서 소수점 셋째자리에서 반올림하여 둘째자리까지 표시함 \par
2. 지가변동률 산정시 백분율로서 소수점 넷째자리에서 반올림하여 셋째자리까지 표시함 \par
3. 지역요인비교치 및 개별요인비교치는 소수점 넷째자리에서 반올림하여 셋째자리까지 표시함 \par
4. 그 밖의 요인 보정치는 표준지 기준 산정방식을 적용할 것 \par
5. 대상토지의 결정단가는 백원단위에서 반올림하여 천원단위까지 표시함 \par
6. 대상건물의 재조달원가 및 감가수정액 등은 백원단위에서 절사하여 천원단위까지 표시함 \par
\end{longquestion}
